\chapter{Applicazioni web con Node.js}

Node.js ci permette di realizzare in modo semplice ed efficiente un server HTTP,
sfruttando i moduli integrati e il suo event loop a thread singolo. Node.js fa
anche un po' di pre-elaborazione sulle richieste (analizza gli header) e ci
fornisce un'astrazione per richieste e risposte.

A parte questo, però, siamo da soli. Per costruire una vera applicazione
dobbiamo affrontare tre problemi:

\begin{itemize}
  \item il \term{routing};
  \item il \term{templating};
  \item l'\term{analisi del corpo} delle richieste.
\end{itemize}

Li affrontiamo uno alla volta, costruendo la stessa applicazione
\guillemotleft lista di cose da fare\guillemotright\ già vista nell'esempio CGI.

\section{Lo scheletro dell'applicazione}

\begin{lstlisting}[style=js, name=app.js]
import http from 'http';

const PORT = 3000;

let server = http.createServer();
server.listen(PORT);
server.on('request', handleRequest);

function handleRequest(request, response) {
  /* gestione della richiesta */
}
\end{lstlisting}

Rispetto all'esempio del capitolo precedente si nota una differenza stilistica:
invece di passare la callback direttamente a \code{createServer}, la si registra
come gestore dell'evento \code{request}. È lo stesso meccanismo degli eventi
del browser visto nel capitolo 7.

\section{Routing}

La nostra applicazione dovrà gestire molti tipi diversi di richiesta HTTP:
mostrare la homepage, mostrare la lista, azzerare l'applicazione, e così via.

\dfn{Routing}{
  Il \term{routing} è il problema di stabilire, data una richiesta, \textbf{quale
  codice} debba essere eseguito per servirla.}

\begin{figure}[H]
  \centering
  \begin{tikzpicture}[font=\Lato\small,
      n/.style={wtnode, minimum width=3.0cm, minimum height=1.0cm},
      a/.style={-{Stealth[length=5pt,width=4pt]}, line width=0.85pt,
                draw=neutral!45!black},
      c/.style={font=\FiraCodeNL\scriptsize, anchor=east,
                text=neutral!25!black, align=right}]

    \node[c] (r1) at (-0.3,0.8)  {GET /todos HTTP/1.1};
    \node[c] (r2) at (-0.3,-0.8) {POST /todos HTTP/1.1\\ todo=Learn+Node.js};

    \node[n, draw=orange-gradient-6, fill=orange-gradient-1!40!white]
      (rt) at (2.9,0) {\textbf{Router}};

    \node[wtnodeplain, anchor=west, font=\FiraCodeNL\scriptsize]
      (h1) at (5.3,0.8) {handleGetReqs()};
    \node[wtnodeplain, anchor=west, font=\FiraCodeNL\scriptsize]
      (h2) at (5.3,-0.8) {handlePostReqs()};

    \draw[a] (r1.east) -- (rt.west |- r1);
    \draw[a] (r2.east) -- (rt.west |- r2);
    \draw[a] (rt.east |- h1) -- (h1.west);
    \draw[a] (rt.east |- h2) -- (h2.west);
  \end{tikzpicture}
  \caption{Il router smista ogni richiesta verso il codice che la serve.}
  \label{fig:router}
\end{figure}

\subsection{Instradare in base al percorso}

\begin{lstlisting}[style=js, name=router.js]
function handleRequest(request, response) {
  switch (request.url) {
    case "/":
      renderHomepage(request, response);
      break;
    case "/reset":
      handleResetRequest(request, response);
      break;
    case "/todo":
      handleTodoListRequest(request, response);
      break;
    /* e cosi' via... */
    default:
      handleError(request, response, 404);
  }
}
\end{lstlisting}

\subsection{Instradare anche in base al metodo}

Le richieste verso un certo percorso possono essere gestite in modo diverso a
seconda del \textbf{metodo HTTP} usato:

\begin{itemize}
  \item una \code{GET} su \code{/todos} restituisce la lista degli elementi;
  \item una \code{POST} su \code{/todos} prova a salvarne uno nuovo.
\end{itemize}

\begin{lstlisting}[style=js, name=router.js]
function handleTodoListRequest(request, response, context = {}) {
  switch (request.method) {
    case "GET":
      handleTodoListRequestGet(request, response, context); break;
    case "POST":
      handleTodoListRequestPost(request, response, context); break;
    default:
      handleError(request, response, 405, "Metodo non supportato");
  }
}
\end{lstlisting}

\info{Il codice 405}{
  Quando il percorso esiste ma il metodo non è previsto, la risposta corretta
  non è un \code{404} (\textit{Not Found}) bensì un \code{405} (\textit{Method
  Not Allowed}). È una distinzione che i client e i \textit{proxy} sanno
  interpretare, e che torna utile soprattutto nelle API REST.}

\subsection{Servire i file statici}

Se il nostro output HTML include file statici (immagini, CSS, script), dobbiamo
configurare il router perché li serva.

\begin{lstlisting}[style=js, name=router.js]
switch (request.url) {
  /* ... */
  case "/css/bootstrap.css":
    fs.readFile('./static/css/bootstrap.css', function (err, data) {
      response.writeHead(200, { 'Content-Type': 'text/css' });
      response.end(data, 'utf-8');
    });
    break;

  case "/img/node.png":
    fs.readFile('./static/img/node.png', function (err, data) {
      response.writeHead(200, { 'Content-Type': 'image/png' });
      response.end(data, "binary");
    });
    break;
}
\end{lstlisting}

Si noti che \code{fs.readFile} è \textbf{asincrona}: prende una callback che
verrà invocata quando il file sarà stato letto. È esattamente il modello non
bloccante descritto nel capitolo precedente.

\subsection{Gestire gli errori}

Quando nessuna delle rotte note si applica, restituiamo un \code{404}.

\begin{lstlisting}[style=js, name=router.js]
switch (request.url) {
  /* altre rotte */
  default:
    handleError(request, response, 404, "Pagina non trovata");
}

function handleError(request, response, statusCode, message) {
  let renderedContent = pug.renderFile("./templates/errorPage.pug",
    { "code": statusCode, "description": message });
  response.writeHead(statusCode, { "Content-Type": "text/html" });
  response.end(renderedContent);
}
\end{lstlisting}

\warning{Il routing non resta semplice a lungo}{
  Uno \code{switch} sembra sufficiente, ma il problema cresce in fretta con
  l'applicazione:
  \begin{itemize}
    \item alcune rotte devono essere accessibili solo ad alcuni utenti (per
          esempio gli amministratori);
    \item alcune rotte dipendono da \textbf{schemi o parametri} nel percorso,
          come \code{GET /users/42/friends}, dove \code{42} può essere un id
          qualunque: con uno \code{switch} su una stringa esatta è impossibile;
    \item può servire eseguire \textbf{due o più pezzi di codice} per gestire una
          richiesta, per esempio una funzione che registra i dati e un'altra che
          prepara la risposta.
  \end{itemize}
  Sono esattamente i problemi che risolveremo con Express, nel prossimo
  capitolo.}

\section{Motori di template}

A un certo punto la nostra applicazione dovrà produrre del codice HTML da
rimandare nel corpo della risposta. Nel primo esempio l'abbiamo scritto a mano,
concatenando stringhe.

Il problema è che alcune parti dell'HTML si ripetono in più pagine (barra di
navigazione, piè di pagina) e che, man mano che le pagine diventano complesse,
costruire una risposta manipolando e concatenando stringhe diventa rapidamente
impraticabile.

\dfn{Motore di template}{
  Un \term{motore di template} (\textit{template engine}) permette di descrivere
  la struttura di una pagina in un file separato, con dei
  \textbf{segnaposto}, e di generare l'HTML finale combinando quel modello con i
  dati forniti a tempo di esecuzione.}

\subsection{Pug}

\term{Pug} (già noto come Jade) è un motore di template ad alte prestazioni
implementato in JavaScript, con una sintassi \textbf{sensibile agli spazi}
(simile a Python) per scrivere modelli HTML.

\begin{lstlisting}[style=shell]
npm install pug
\end{lstlisting}

\begin{lstlisting}[style=js]
import pug from "pug";

let html = pug.renderFile("./templates/basic_example.pug");
console.log(html);   // html contiene il codice HTML generato
\end{lstlisting}

\begin{esempio}{Il primo template Pug}
\begin{lstlisting}[style=pug, name=basic_example.pug]
doctype html
html(lang="it")
  head
    title Hello Pug!
  body
    h1(class="foo") Primo template Pug
    p Questo e' il nostro primo template Pug!
\end{lstlisting}

produce:

\begin{lstlisting}[style=html, name=output.html]
<!DOCTYPE html>
<html lang="it">
  <head><title>Hello Pug!</title></head>
  <body>
    <h1 class="foo">Primo template Pug</h1>
    <p>Questo e' il nostro primo template Pug!</p>
  </body>
</html>
\end{lstlisting}

L'indentazione sostituisce i tag di chiusura, e gli attributi si scrivono fra
parentesi tonde. Il risultato è molto più breve e, soprattutto, impossibile da
sbagliare nell'annidamento.
\end{esempio}

\subsection{Inclusioni}

Pug è molto più di un modo diverso di scrivere HTML. Tanto per cominciare, un
template può \textbf{includerne} altri: è un grande aiuto per il riuso.

\begin{lstlisting}[style=pug, name=partials/head.pug]
head
  title Hello Pug
  script(src='/js/script.js')
  link(rel='stylesheet' href='/style.css')
\end{lstlisting}

\begin{lstlisting}[style=pug, name=partials/footer.pug]
footer#footer
  p Copyright (c) Web Tech
\end{lstlisting}

\begin{lstlisting}[style=pug, name=index.pug]
doctype html
html
  include partials/head.pug
  body
    h1 Hello Pug!
    include partials/footer.pug
\end{lstlisting}

Si noti \code{footer\#footer}: in Pug il cancelletto è l'abbreviazione per
l'attributo \attr{id}, esattamente come nei selettori CSS.

\subsection{Ereditarietà fra template}

Progettando applicazioni web è comune avere un \textbf{layout condiviso} da
tutte le pagine, che ciascuna pagina può poi personalizzare in alcune parti. Pug
supporta questo scenario con l'\term{ereditarietà fra template}, tramite le
parole chiave \code{block} ed \code{extends}.

\begin{figure}[H]
  \centering
  \begin{tikzpicture}[font=\Lato\scriptsize,
      b/.style={draw=primary!60!black, fill=primary!10!white, line width=0.7pt,
                minimum width=3.0cm, minimum height=0.5cm, anchor=north},
      bc/.style={b, draw=accent!60!black, fill=accent!12!white}]

    \foreach \x/\ttl in {0/{Layout base}, 4.4/{Homepage}, 8.8/{Pagina articolo}} {
      \node[font=\Lato\scriptsize\bfseries, text=neutral!25!black]
        at (\x,1.0) {\ttl};
    }

    % layout base
    \node[b] (h0) at (0,0.6) {heading};
    \node[b, below=1pt of h0] (n0) {navbar};
    \node[bc, below=1pt of n0, minimum height=1.1cm] (c0) {content};
    \node[b, below=1pt of c0] (f0) {footer};

    % homepage
    \node[b] (h1) at (4.4,0.6) {heading};
    \node[b, below=1pt of h1] (n1) {navbar};
    \node[bc, below=1pt of n1, minimum height=1.1cm] (c1) {contenuto proprio};
    \node[b, below=1pt of c1] (f1) {footer};

    % articolo
    \node[b] (h2) at (8.8,0.6) {heading};
    \node[b, below=1pt of h2] (n2) {navbar};
    \node[bc, below=1pt of n2, minimum height=1.1cm, minimum width=1.9cm,
          anchor=north west] (c2) at ([xshift=-1.5cm]n2.south west) {article};
    \node[bc, minimum height=1.1cm, minimum width=1.0cm, anchor=north west]
      (s2) at ([xshift=1pt]c2.north east) {sidebar};
    \node[b, below=1pt of c2.south west, anchor=north west] (f2) {footer};
  \end{tikzpicture}
  \caption{Un layout comune, personalizzato dalle singole pagine.}
  \label{fig:pug-inheritance}
\end{figure}

\begin{esempio}{Layout base e pagina figlia}
\begin{lstlisting}[style=pug, name=base-layout.pug]
doctype html
html
  head
    title Pug Inheritance
  body
    block heading
      h1 Pug Inheritance
    block content
    block footer
      footer Questo e' il footer predefinito.
\end{lstlisting}

Il template definisce tre \textbf{blocchi}: \code{heading}, \code{content} e
\code{footer}. I blocchi \code{heading} e \code{footer} contengono un contenuto
predefinito.

\begin{lstlisting}[style=pug, name=index.pug]
extends ./base-layout.pug

block content
  p Il contenuto vero e proprio della homepage.

block footer
  footer Questo e' un footer specializzato.
\end{lstlisting}

Il template figlio \textbf{sovrascrive} i blocchi \code{content} e
\code{footer}; il blocco \code{heading} non viene sovrascritto e mostra il
contenuto predefinito. Il risultato è:

\begin{lstlisting}[style=html, name=output.html]
<!DOCTYPE html>
<html>
  <head><title>Pug Inheritance</title></head>
  <body>
    <h1>Pug Inheritance</h1>
    <p>Il contenuto vero e proprio della homepage.</p>
    <footer>Questo e' un footer specializzato.</footer>
  </body>
</html>
\end{lstlisting}
\end{esempio}

\subsection{Locals: passare i dati al template}

I template Pug possono generare contenuto a partire da un oggetto di dati
fornito, chiamato \term{locals}, passato alle funzioni \code{render()} o
\code{renderFile()}.

\begin{esempio}{Interpolazione}
\begin{lstlisting}[style=js]
import pug from "pug";

let html = pug.renderFile("./locals.pug", {
  "product": {
    "name": "Samsung S24",
    "description": "Smartphone"
  }
});

console.log(html);
\end{lstlisting}

\begin{lstlisting}[style=pug, name=locals.pug]
- let footer = "una bella stringa"

h1 #{product.name.toUpperCase()}
p Descrizione: #{product.description}
footer #{footer}
\end{lstlisting}

\begin{lstlisting}[style=html, name=output.html]
<h1>SAMSUNG S24</h1>
<p>Descrizione: Smartphone</p>
<footer>una bella stringa</footer>
\end{lstlisting}

Il codice contenuto fra \code{\#\{} e \code{\}} viene valutato,
\textbf{sottoposto a escaping} e inserito nell'output. L'escaping automatico
non è un dettaglio: è ciò che impedisce che un nome di prodotto contenente
\code{<script>} finisca per essere eseguito dal browser.

Una riga che inizia con \code{-} contiene invece codice JavaScript eseguito ma
non stampato.
\end{esempio}

\subsection{Condizionali}

\begin{lstlisting}[style=js]
let user = { "isAdmin": true, "name": "Brendan" };
let html = pug.renderFile("./conditionals.pug", { "user": user });
\end{lstlisting}

\begin{lstlisting}[style=pug, name=conditionals.pug]
if !user
  a(href="/login") Per favore, autenticati.
else if user.isAdmin
  p Ai tuoi ordini, #{user.name}
else
  p Bentornato, #{user.name}
\end{lstlisting}

\subsection{Iterazioni}

\begin{lstlisting}[style=js]
let items = [
  { "name": "Margherita",  "price": 5.50 },
  { "name": "Marinara",    "price": 5.00 },
  { "name": "Capricciosa", "price": 6.50 }
];
let html = pug.renderFile("./iteration.pug", { "items": items });
\end{lstlisting}

\begin{lstlisting}[style=pug, name=iteration.pug]
ul
  each item in items
    li #{item.name} - EUR #{item.price.toFixed(2)}
  else
    li Nessuna pizza disponibile al momento!
\end{lstlisting}

\begin{lstlisting}[style=html, name=output.html]
<ul>
  <li>Margherita - EUR 5.50</li>
  <li>Marinara - EUR 5.00</li>
  <li>Capricciosa - EUR 6.50</li>
</ul>
\end{lstlisting}

Il ramo \code{else} di un \code{each} viene eseguito quando la collezione è
vuota: è una comodità che evita di scrivere a mano il controllo.

\info{Non solo Pug}{
  Pug è solo uno dei tanti motori di template disponibili. Altri molto noti sono
  \term{EJS} (con una sintassi simile a JSP e PHP), \term{SquirrellyJS},
  \term{Nunjucks}, \term{LiquidJS}, \term{Eta} e \term{Haml} (per Ruby). Cambia
  la sintassi, non il concetto.}

\section{Analizzare il corpo delle richieste}

Gli utenti compilano il form della pagina per salvare nuovi elementi. Il form
viene inviato con il metodo \code{POST}, quindi dobbiamo analizzare il corpo
della richiesta per ottenere i nomi dei parametri e i loro valori. Non è così
semplice come sembra.

\begin{lstlisting}[style=pug, name=todo.pug]
form(action="/todo" method="POST")
  label(for="todo") Elemento da fare
  input(type="text" id="todo" name="todo" required)
  button(type="submit") Salva
\end{lstlisting}

\begin{lstlisting}[style=http, name=Richiesta generata]
POST /todo HTTP/1.1

todo=Learn+Node.js
\end{lstlisting}

\subsection{Il corpo arriva a pezzi}

Il corpo della richiesta \textbf{potrebbe non essere ancora disponibile} quando
iniziamo a elaborarla: magari l'utente ha una connessione lenta, o sta caricando
file di grandi dimensioni, e il corpo diventerà disponibile più tardi.

L'oggetto \code{request} (di tipo \code{http.IncomingMessage}) estende
\code{stream.Readable}: genera un evento \code{data} ogni volta che un nuovo
pezzo di dati diventa disponibile, e un evento \code{end} quando non c'è più
nulla da consumare. Per leggere il corpo dobbiamo quindi operare in modo
\textbf{asincrono}, sfruttando questi due eventi.

\begin{lstlisting}[style=js, name=body-parser.js]
function parseRequestBody(request) {
  return new Promise((resolve, reject) => {
    let body = [];
    request
      .on('data', chunk => { body.push(chunk); })
      .on('end', () => {
        body = Buffer.concat(body).toString();
        // qui `body` contiene l'intero corpo come stringa
        let data = parseRequestBodyString(body);
        resolve(data);
      });
  });
}
\end{lstlisting}

Si notino i due meccanismi che abbiamo studiato separatamente e che qui si
incontrano: gli \textbf{eventi} (capitolo 7) per raccogliere i pezzi, e una
\textbf{Promise} (capitolo 8) per esporre il risultato a chi ha chiamato.

\subsection{Interpretare la stringa}

Letto il corpo come stringa, si tratta di suddividerla per ottenere nomi e
valori dei parametri. Ricordiamo che il corpo è una stringa della forma
\code{p1=v1\&p2=v2\&...}, e che è \textbf{codificato come URL}.

\begin{lstlisting}[style=js, name=body-parser.js]
function parseRequestBodyString(body) {
  let data = {};
  let slices = body.split("&");

  for (let slice of slices) {
    let paramName = slice.split("=")[0];
    paramName = decodeURIComponent(paramName.replace(/\+/g, " "));

    let paramValue = slice.split("=")[1];
    paramValue = decodeURIComponent(paramValue.replace(/\+/g, " "));

    data[paramName] = paramValue;
  }
  return data;
}
\end{lstlisting}

\warning{Ricordarsi della codifica}{
  Il corpo è codificato come URL, esattamente come descritto nel capitolo su
  HTML: la frase \guillemotleft Prove that P=NP\guillemotright\ viaggia come
  \code{Prove+that+P\%3DNP}. Bisogna quindi sostituire i \code{+} con spazi
  \textbf{e} applicare \code{decodeURIComponent}: saltare uno dei due passaggi
  produce bug che si manifestano solo con certi input.}

\section{Anatomia di un'applicazione web}

Mettendo insieme i pezzi visti finora, si delinea l'architettura tipica di
un'applicazione web.

\begin{figure}[H]
  \centering
  \begin{tikzpicture}[font=\Lato\small,
      n/.style={wtnode, minimum width=3.0cm, minimum height=1.0cm},
      a/.style={-{Stealth[length=5pt,width=4pt]}, line width=0.85pt,
                draw=neutral!45!black}]

    \node[n, minimum height=1.8cm] (br) at (0,0)
      {\textbf{Web Browser}\\[3pt]\scriptsize HTML, CSS, JS};

    \node[n, draw=orange-gradient-6, fill=orange-gradient-1!40!white]
      (rt) at (4.6,1.3) {\textbf{Router}};
    \node[n] (bl) at (4.6,0) {\textbf{Logica di business}\\\scriptsize (controller)};
    \node[n] (pr) at (4.6,-1.4) {\textbf{Presentazione}\\\scriptsize (template)};
    \node[n, draw=green-gradient-6, fill=green-gradient-1!45!white]
      (db) at (8.8,0) {\textbf{Persistenza}\\\scriptsize dei dati};

    \draw[a] (br.east) to[out=25, in=180] node[wtlab, above, pos=0.5]
      {richiesta HTTP} (rt.west);
    \draw[a] (rt) -- (bl);
    \draw[a] (bl) -- (pr);
    \draw[a] (bl) -- (db);
    \draw[a, draw=accent!70!black] (pr.west) to[out=180, in=-25]
      node[wtlab, below, pos=0.5]{risposta HTTP} (br.east);

    \begin{scope}[on background layer]
      \node[rounded corners=5pt, draw=primary!40!white, line width=0.9pt,
            fit=(rt)(bl)(pr)(db), inner sep=11pt,
            label={[font=\Lato\scriptsize, text=primary!60!black]above:Backend}] {};
      \node[rounded corners=5pt, draw=primary!40!white, line width=0.9pt,
            fit=(br), inner sep=11pt,
            label={[font=\Lato\scriptsize, text=primary!60!black]above:Frontend}] {};
    \end{scope}
  \end{tikzpicture}
  \caption{L'architettura tipica di un'applicazione web tradizionale.}
  \label{fig:webapp-anatomy}
\end{figure}

\section{Riferimenti}

\begin{itemize}
  \item \textbf{Web Templating} --- EduTech wiki, Università di Ginevra. \\
        \urlx{https://edutechwiki.unige.ch/en/Web_templating}
  \item \textbf{Single page apps in depth}, Mikito Takada. \\
        \urlx{http://singlepageappbook.com/} \\
        Parte rilevante: \textit{Templating: from data to HTML}.
  \item \textbf{Node.js v.20.X (LTS) documentation}. \\
        \urlx{https://nodejs.org/docs/latest-v20.x/api/index.html}
  \item \textbf{A beginner's Guide to Pug}, James Hibbard. \\
        \urlx{https://www.sitepoint.com/a-beginners-guide-to-pug/}
\end{itemize}
